\iftestbuild % TEST-ONLY-BEGIN
\setcounter{section}{2}\setcounter{theorem}{28}
\fi % TEST-ONLY-END
% Source: book pp. 69--81 (PDF 83--95); solutions pp. 83--92.
\section{Matrices}

\subsection{Representing a Linear Map by a Matrix}

We know that if $v_1,\dots,v_n$ is a basis of $V$ and $T\colon V\to W$ is
linear, then the values of $Tv_1,\dots,Tv_n$ determine the values of $T$ on
arbitrary vectors in $V$---see the linear map lemma (\ref{ax:3.4}). As we will
soon see, matrices provide an efficient method of recording the values of the
$Tv_k$'s in terms of a basis of $W$.

\begin{definition}[Matrix, $A_{j,k}$]\label{ax:3.29}
Suppose $m$ and $n$ are nonnegative integers. An \emph{$m$-by-$n$ matrix} $A$
is a rectangular array of elements of $\F$ with $m$ rows and $n$ columns:
\[ A = \begin{pmatrix} A_{1,1} & \cdots & A_{1,n}\\ \vdots & & \vdots\\
   A_{m,1} & \cdots & A_{m,n} \end{pmatrix}. \]
The notation $A_{j,k}$ denotes the entry in row $j$, column $k$ of $A$.
\end{definition}

\marginnote{When dealing with matrices, the first index refers to the row
number; the second index refers to the column number.}%
\begin{example}[$A_{j,k}$ equals entry in row $j$, column $k$ of $A$]\label{ax:3.30}
Suppose $A = \begin{pmatrix} 8 & 4 & 5-3i\\ 1 & 9 & 7 \end{pmatrix}$.
Thus $A_{2,3}$ refers to the entry in the second row, third column of $A$,
which means that $A_{2,3}=7$.
\end{example}

Now we come to the key definition in this section.

\begin{definition}[Matrix of a linear map, $\Mat(T)$]\label{ax:3.31}
Suppose $T\in\Lin(V,W)$ and $v_1,\dots,v_n$ is a basis of $V$ and
$w_1,\dots,w_m$ is a basis of $W$. The \emph{matrix of $T$} with respect to
these bases is the $m$-by-$n$ matrix $\Mat(T)$ whose entries $A_{j,k}$ are
defined by
\[ Tv_k = A_{1,k}w_1+\cdots+A_{m,k}w_m. \]
If the bases $v_1,\dots,v_n$ and $w_1,\dots,w_m$ are not clear from the
context, then the notation $\Mat\bigl(T,(v_1,\dots,v_n),(w_1,\dots,w_m)\bigr)$
is used.
\end{definition}

The matrix $\Mat(T)$ of a linear map $T\in\Lin(V,W)$ depends on the basis
$v_1,\dots,v_n$ of $V$ and the basis $w_1,\dots,w_m$ of $W$, as well as on $T$.
However, the bases should be clear from the context, and thus they are often
not included in the notation.

To remember how $\Mat(T)$ is constructed from $T$, you might write across the
top of the matrix the basis vectors $v_1,\dots,v_n$ for the domain and along
the left the basis vectors $w_1,\dots,w_m$ for the vector space into which $T$
maps, as follows:
\[ \Mat(T) = \bordermatrix{
      & v_1 & \cdots & v_k & \cdots & v_n \cr
  w_1 &     &        & A_{1,k} &    &     \cr
  \vdots &  &        & \vdots  &    &     \cr
  w_m &     &        & A_{m,k} &    &     \cr}. \]

\marginnote{The $k^{\text{th}}$ column of $\Mat(T)$ consists of the scalars
needed to write $Tv_k$ as a linear combination of $w_1,\dots,w_m$:
\[ Tv_k = \sum_{j=1}^m A_{j,k}w_j. \]}%
In the matrix above only the $k^{\text{th}}$ column is shown. Thus the second
index of each displayed entry of the matrix above is $k$. The picture above
should remind you that $Tv_k$ can be computed from $\Mat(T)$ by multiplying
each entry in the $k^{\text{th}}$ column by the corresponding $w_j$ from the
left column, and then adding up the resulting vectors.

\marginnote{If $T$ is a linear map from an $n$-dimensional vector space to an
$m$-dimensional vector space, then $\Mat(T)$ is an $m$-by-$n$ matrix.}%
If $T$ is a linear map from $\F^n$ to $\F^m$, then unless stated otherwise,
assume the bases in question are the standard ones (where the
$k^{\text{th}}$ basis vector is $1$ in the $k^{\text{th}}$ slot and $0$ in all
other slots). If you think of elements of $\F^m$ as columns of $m$ numbers,
then you can think of the $k^{\text{th}}$ column of $\Mat(T)$ as $T$ applied
to the $k^{\text{th}}$ standard basis vector.

\begin{example}[The matrix of a linear map from $\F^2$ to $\F^3$]\label{ax:3.32}
Suppose $T\in\Lin(\F^2,\F^3)$ is defined by
\[ T(x,y) = (x+3y,\,2x+5y,\,7x+9y). \]
Because $T(1,0)=(1,2,7)$ and $T(0,1)=(3,5,9)$, the matrix of $T$ with respect
to the standard bases is the $3$-by-$2$ matrix below:
\[ \Mat(T) = \begin{pmatrix} 1 & 3\\ 2 & 5\\ 7 & 9 \end{pmatrix}. \]
\end{example}

When working with $\Poly_m(\F)$, use the standard basis
$1,x,x^2,\dots,x^m$ unless the context indicates otherwise.

\begin{example}[Matrix of the differentiation map from $\Poly_3(\R)$ to
$\Poly_2(\R)$]\label{ax:3.33}
Suppose $D\in\Lin\bigl(\Poly_3(\R),\allowbreak\Poly_2(\R)\bigr)$ is the differentiation
map defined by $Dp=p'$. Because $(x^n)'=nx^{n-1}$, the matrix of $D$ with
respect to the standard bases is the $3$-by-$4$ matrix below:
\[ \Mat(D) = \begin{pmatrix} 0 & 1 & 0 & 0\\ 0 & 0 & 2 & 0\\ 0 & 0 & 0 & 3
   \end{pmatrix}. \]
\end{example}

\subsection{Addition and Scalar Multiplication of Matrices}

For the rest of this section, assume that $U$, $V$, and $W$ are
finite-dimensional and that a basis has been chosen for each of these vector
spaces. Thus for each linear map from $V$ to $W$, we can talk about its matrix
(with respect to the chosen bases).

Is the matrix of the sum of two linear maps equal to the sum of the matrices of
the two maps? Right now this question does not yet make sense because although
we have defined the sum of two linear maps, we have not defined the sum of two
matrices. Fortunately, the natural definition of the sum of two matrices has
the right properties. Specifically, we make the following definition.

\begin{definition}[Matrix addition]\label{ax:3.34}
The \emph{sum of two matrices of the same size} is the matrix obtained by
adding corresponding entries in the matrices:
\begin{multline*}
\begin{pmatrix} A_{1,1} & \cdots & A_{1,n}\\ \vdots & & \vdots\\
   A_{m,1} & \cdots & A_{m,n} \end{pmatrix}
 + \begin{pmatrix} C_{1,1} & \cdots & C_{1,n}\\ \vdots & & \vdots\\
   C_{m,1} & \cdots & C_{m,n} \end{pmatrix}\\
 = \begin{pmatrix} A_{1,1}+C_{1,1} & \cdots & A_{1,n}+C_{1,n}\\
   \vdots & & \vdots\\
   A_{m,1}+C_{m,1} & \cdots & A_{m,n}+C_{m,n} \end{pmatrix}.
\end{multline*}
\end{definition}

In the next result, the assumption is that the same bases are used for all
three linear maps $S+T$, $S$, and $T$.

\begin{result}[Matrix of the sum of linear maps]\label{ax:3.35}
Suppose $S,T\in\Lin(V,W)$. Then $\Mat(S+T)=\Mat(S)+\Mat(T)$.
\end{result}

The verification of the result above follows from the definitions and is left
to the reader.

Still assuming that we have some bases in mind, is the matrix of a scalar times
a linear map equal to the scalar times the matrix of the linear map? Again, the
question does not yet make sense because we have not defined scalar
multiplication on matrices. Fortunately, the natural definition again has the
right properties.

\begin{definition}[Scalar multiplication of a matrix]\label{ax:3.36}
The product of a scalar and a matrix is the matrix obtained by multiplying each
entry in the matrix by the scalar:
\[ \lambda\begin{pmatrix} A_{1,1} & \cdots & A_{1,n}\\ \vdots & & \vdots\\
   A_{m,1} & \cdots & A_{m,n} \end{pmatrix}
 = \begin{pmatrix} \lambda A_{1,1} & \cdots & \lambda A_{1,n}\\
   \vdots & & \vdots\\
   \lambda A_{m,1} & \cdots & \lambda A_{m,n} \end{pmatrix}. \]
\end{definition}

\begin{example}[Addition and scalar multiplication of matrices]\label{ax:3.37}
\leavevmode
\[ 2\begin{pmatrix} 3 & 1\\ -1 & 5 \end{pmatrix}
   + \begin{pmatrix} 4 & 2\\ 1 & 6 \end{pmatrix}
 = \begin{pmatrix} 6 & 2\\ -2 & 10 \end{pmatrix}
   + \begin{pmatrix} 4 & 2\\ 1 & 6 \end{pmatrix}
 = \begin{pmatrix} 10 & 4\\ -1 & 16 \end{pmatrix} \]
\end{example}

In the next result, the assumption is that the same bases are used for both the
linear maps $\lambda T$ and $T$.

\begin{result}[The matrix of a scalar times a linear map]\label{ax:3.38}
Suppose $\lambda\in\F$ and $T\in\Lin(V,W)$. Then
$\Mat(\lambda T)=\lambda\Mat(T)$.
\end{result}

The verification of the result above is also left to the reader.

Because addition and scalar multiplication have now been defined for matrices,
you should not be surprised that a vector space is about to appear. First we
introduce a bit of notation so that this new vector space has a name, and then
we find the dimension of this new vector space.

\begin{notation}[$\F^{m,n}$]\label{ax:3.39}
For $m$ and $n$ positive integers, the set of all $m$-by-$n$ matrices with
entries in $\F$ is denoted by $\F^{m,n}$.
\end{notation}

\begin{result}[$\dim\F^{m,n}=mn$]\label{ax:3.40}
Suppose $m$ and $n$ are positive integers. With addition and scalar
multiplication defined as above, $\F^{m,n}$ is a vector space of dimension
$mn$.
\end{result}

\begin{proof}
The verification that $\F^{m,n}$ is a vector space is left to the reader. Note
that the additive identity of $\F^{m,n}$ is the $m$-by-$n$ matrix all of whose
entries equal $0$.

The reader should also verify that the list of distinct $m$-by-$n$ matrices
that have $0$ in all entries except for a $1$ in one entry is a basis of
$\F^{m,n}$. There are $mn$ such matrices, so the dimension of $\F^{m,n}$ equals
$mn$.
\end{proof}

\subsection{Matrix Multiplication}

Suppose, as previously, that $v_1,\dots,v_n$ is a basis of $V$ and
$w_1,\dots,w_m$ is a basis of $W$. Suppose also that $u_1,\dots,u_p$ is a basis
of $U$.

Consider linear maps $T\colon U\to V$ and $S\colon V\to W$. The composition
$ST$ is a linear map from $U$ to $W$. Does $\Mat(ST)$ equal $\Mat(S)\Mat(T)$?
This question does not yet make sense because we have not defined the product
of two matrices. We will choose a definition of matrix multiplication that
forces this question to have a positive answer. Let's see how to do this.

Suppose $\Mat(S)=A$ and $\Mat(T)=B$. For $1\le k\le p$, we have
\begin{align*}
(ST)u_k &= S\Bigl(\sum_{r=1}^n B_{r,k}v_r\Bigr)\\
        &= \sum_{r=1}^n B_{r,k}Sv_r\\
        &= \sum_{r=1}^n B_{r,k}\sum_{j=1}^m A_{j,r}w_j\\
        &= \sum_{j=1}^m\Bigl(\sum_{r=1}^n A_{j,r}B_{r,k}\Bigr)w_j.
\end{align*}
Thus $\Mat(ST)$ is the $m$-by-$p$ matrix whose entry in row $j$, column $k$,
equals
\[ \sum_{r=1}^n A_{j,r}B_{r,k}. \]

Now we see how to define matrix multiplication so that the desired equation
$\Mat(ST)=\Mat(S)\Mat(T)$ holds.

\begin{definition}[Matrix multiplication]\label{ax:3.41}
Suppose $A$ is an $m$-by-$n$ matrix and $B$ is an $n$-by-$p$ matrix. Then $AB$
is defined to be the $m$-by-$p$ matrix whose entry in row $j$, column $k$, is
given by the equation
\[ (AB)_{j,k} = \sum_{r=1}^n A_{j,r}B_{r,k}. \]
Thus the entry in row $j$, column $k$, of $AB$ is computed by taking row $j$ of
$A$ and column $k$ of $B$, multiplying together corresponding entries, and then
summing.
\end{definition}

\marginnote{You may have learned this definition of matrix multiplication in
an earlier course, although you may not have seen this motivation for it.}%
Note that we define the product of two matrices only when the number of columns
of the first matrix equals the number of rows of the second matrix.

\begin{example}[Matrix multiplication]\label{ax:3.42}
Here we multiply together a $3$-by-$2$ matrix and a $2$-by-$4$ matrix,
obtaining a $3$-by-$4$ matrix:
\[ \begin{pmatrix} 1 & 2\\ 3 & 4\\ 5 & 6 \end{pmatrix}
   \begin{pmatrix} 6 & 5 & 4 & 3\\ 2 & 1 & 0 & -1 \end{pmatrix}
 = \begin{pmatrix} 10 & 7 & 4 & 1\\ 26 & 19 & 12 & 5\\ 42 & 31 & 20 & 9
   \end{pmatrix}. \]
\end{example}

Matrix multiplication is not commutative---$AB$ is not necessarily equal to
$BA$ even if both products are defined (see Exercise~10). Matrix multiplication
is distributive and associative (see Exercises 11 and~12).

In the next result, we assume that the same basis of $V$ is used in considering
$T\in\Lin(U,V)$ and $S\in\Lin(V,W)$, the same basis of $W$ is used in
considering $S\in\Lin(V,W)$ and $ST\in\Lin(U,W)$, and the same basis of $U$ is
used in considering $T\in\Lin(U,V)$ and $ST\in\Lin(U,W)$.

\begin{result}[Matrix of product of linear maps]\label{ax:3.43}
If $T\in\Lin(U,V)$ and $S\in\Lin(V,W)$, then $\Mat(ST)=\Mat(S)\Mat(T)$.
\end{result}

The proof of the result above is the calculation that was done as motivation
before the definition of matrix multiplication.

In the next piece of notation, note that as usual the first index refers to a
row and the second index refers to a column, with a vertically centered dot
used as a placeholder.

\begin{notation}[$A_{j,\cdot}$, $A_{\cdot,k}$]\label{ax:3.44}
Suppose $A$ is an $m$-by-$n$ matrix.
\begin{itemize}
\item If $1\le j\le m$, then $A_{j,\cdot}$ denotes the $1$-by-$n$ matrix
  consisting of row $j$ of $A$.
\item If $1\le k\le n$, then $A_{\cdot,k}$ denotes the $m$-by-$1$ matrix
  consisting of column $k$ of $A$.
\end{itemize}
\end{notation}

\begin{example}[$A_{j,\cdot}$ equals $j^{\text{th}}$ row of $A$ and
$A_{\cdot,k}$ equals $k^{\text{th}}$ column of $A$]\label{ax:3.45}
The notation $A_{2,\cdot}$ denotes the second row of $A$ and $A_{\cdot,2}$
denotes the second column of $A$. Thus if
$A=\begin{pmatrix} 8 & 4 & 5\\ 1 & 9 & 7 \end{pmatrix}$, then
\[ A_{2,\cdot} = \begin{pmatrix} 1 & 9 & 7 \end{pmatrix}
   \quad\text{and}\quad
   A_{\cdot,2} = \begin{pmatrix} 4\\ 9 \end{pmatrix}. \]
\end{example}

The product of a $1$-by-$n$ matrix and an $n$-by-$1$ matrix is a $1$-by-$1$
matrix. However, we will frequently identify a $1$-by-$1$ matrix with its
entry. For example,
\[ \begin{pmatrix} 3 & 4 \end{pmatrix}\begin{pmatrix} 6\\ 2 \end{pmatrix}
   = \begin{pmatrix} 26 \end{pmatrix} \]
because $3\cdot6+4\cdot2=26$. However, we can identify
$\begin{pmatrix} 26 \end{pmatrix}$ with $26$, writing
$\begin{pmatrix} 3 & 4 \end{pmatrix}\begin{pmatrix} 6\\ 2 \end{pmatrix}=26$.

The next result uses the convention discussed in the paragraph above to give
another way to think of matrix multiplication. For example, the next result and
the calculation in the paragraph above explain why the entry in row $2$,
column $1$, of the product in Example~\ref{ax:3.42} equals $26$.

\begin{result}[Entry of matrix product equals row times column]\label{ax:3.46}
Suppose $A$ is an $m$-by-$n$ matrix and $B$ is an $n$-by-$p$ matrix. Then
\[ (AB)_{j,k} = A_{j,\cdot}B_{\cdot,k} \]
if $1\le j\le m$ and $1\le k\le p$. In other words, the entry in row $j$,
column $k$, of $AB$ equals (row $j$ of $A$) times (column $k$ of $B$).
\end{result}

\begin{proof}
Suppose $1\le j\le m$ and $1\le k\le p$. The definition of matrix
multiplication states that
\begin{equation*}
(AB)_{j,k} = A_{j,1}B_{1,k}+\cdots+A_{j,n}B_{n,k}.
\refstepcounter{theorem}\tag*{\thetheorem}\label{ax:3.47}
\end{equation*}
The definition of matrix multiplication also implies that the product of the
$1$-by-$n$ matrix $A_{j,\cdot}$ and the $n$-by-$1$ matrix $B_{\cdot,k}$ is the
$1$-by-$1$ matrix whose entry is the number on the right side of the equation
above. Thus the entry in row $j$, column $k$, of $AB$ equals (row $j$ of $A$)
times (column $k$ of $B$).
\end{proof}

The next result gives yet another way to think of matrix multiplication. In the
result below, $(AB)_{\cdot,k}$ is column $k$ of the $m$-by-$p$ matrix $AB$.
Thus $(AB)_{\cdot,k}$ is an $m$-by-$1$ matrix. Also, $AB_{\cdot,k}$ is an
$m$-by-$1$ matrix because it is the product of an $m$-by-$n$ matrix and an
$n$-by-$1$ matrix. Thus the two sides of the equation in the result below have
the same size, making it reasonable that they might be equal.

\begin{result}[Column of matrix product equals matrix times column]\label{ax:3.48}
Suppose $A$ is an $m$-by-$n$ matrix and $B$ is an $n$-by-$p$ matrix. Then
\[ (AB)_{\cdot,k} = AB_{\cdot,k} \]
if $1\le k\le p$. In other words, column $k$ of $AB$ equals $A$ times column
$k$ of $B$.
\end{result}

\begin{proof}
As discussed above, $(AB)_{\cdot,k}$ and $AB_{\cdot,k}$ are both $m$-by-$1$
matrices. If $1\le j\le m$, then the entry in row $j$ of $(AB)_{\cdot,k}$ is
the left side of \ref{ax:3.47} and the entry in row $j$ of $AB_{\cdot,k}$ is
the right side of \ref{ax:3.47}. Thus $(AB)_{\cdot,k}=AB_{\cdot,k}$.
\end{proof}

Our next result will give another way of thinking about the product of an
$m$-by-$n$ matrix and an $n$-by-$1$ matrix, motivated by the next example.

\begin{example}[Product of a $3$-by-$2$ matrix and a $2$-by-$1$
matrix]\label{ax:3.49}
Use our definitions and basic arithmetic to verify that
\[ \begin{pmatrix} 1 & 2\\ 3 & 4\\ 5 & 6 \end{pmatrix}
   \begin{pmatrix} 5\\ 1 \end{pmatrix}
 = \begin{pmatrix} 7\\ 19\\ 31 \end{pmatrix}
 = 5\begin{pmatrix} 1\\ 3\\ 5 \end{pmatrix}
   + 1\begin{pmatrix} 2\\ 4\\ 6 \end{pmatrix}. \]
Thus in this example, the product of a $3$-by-$2$ matrix and a $2$-by-$1$
matrix is a linear combination of the columns of the $3$-by-$2$ matrix, with
the scalars ($5$ and $1$) that multiply the columns coming from the $2$-by-$1$
matrix.
\end{example}

The next result generalizes the example above.

\begin{result}[Linear combination of columns]\label{ax:3.50}
Suppose $A$ is an $m$-by-$n$ matrix and
$b=\begin{pmatrix} b_1\\ \vdots\\ b_n \end{pmatrix}$ is an $n$-by-$1$ matrix.
Then
\[ Ab = b_1A_{\cdot,1}+\cdots+b_nA_{\cdot,n}. \]
In other words, $Ab$ is a linear combination of the columns of $A$, with the
scalars that multiply the columns coming from $b$.
\end{result}

\begin{proof}
If $k\in\{1,\dots,m\}$, then the definition of matrix multiplication implies
that the entry in row $k$ of the $m$-by-$1$ matrix $Ab$ is
\[ A_{k,1}b_1+\cdots+A_{k,n}b_n. \]
The entry in row $k$ of $b_1A_{\cdot,1}+\cdots+b_nA_{\cdot,n}$ also equals the
number displayed above.

Because $Ab$ and $b_1A_{\cdot,1}+\cdots+b_nA_{\cdot,n}$ have the same entry in
row $k$ for each $k\in\{1,\dots,m\}$, we conclude that
$Ab=b_1A_{\cdot,1}+\cdots+b_nA_{\cdot,n}$.
\end{proof}

Our two previous results focus on the columns of a matrix. Analogous results
hold for the rows of a matrix. Specifically, see Exercises 8 and~9, which can
be proved using appropriate modifications of the proofs of \ref{ax:3.48} and
\ref{ax:3.50}.

The next result is the main tool used in the next subsection to prove the
column--row factorization (\ref{ax:3.56}) and to prove that the column rank of
a matrix equals the row rank (\ref{ax:3.57}). To be consistent with the
notation often used with the column--row factorization, including in the next
subsection, the matrices in the next result are called $C$ and $R$ instead of
$A$ and $B$.

\begin{result}[Matrix multiplication as linear combinations of columns or
rows]\label{ax:3.51}
Suppose $C$ is an $m$-by-$c$ matrix and $R$ is a $c$-by-$n$ matrix.
\begin{parts}
\item If $k\in\{1,\dots,n\}$, then column $k$ of $CR$ is a linear combination
  of the columns of $C$, with the coefficients of this linear combination
  coming from column $k$ of $R$.
\item If $j\in\{1,\dots,m\}$, then row $j$ of $CR$ is a linear combination of
  the rows of $R$, with the coefficients of this linear combination coming
  from row $j$ of $C$.
\end{parts}
\end{result}

\begin{proof}
Suppose $k\in\{1,\dots,n\}$. Then column $k$ of $CR$ equals $CR_{\cdot,k}$ (by
\ref{ax:3.48}), which equals the linear combination of the columns of $C$ with
coefficients coming from $R_{\cdot,k}$ (by \ref{ax:3.50}). Thus (a) holds.

To prove (b), follow the pattern of the proof of (a) but use rows instead of
columns and use Exercises 8 and~9 instead of \ref{ax:3.48} and \ref{ax:3.50}.
\end{proof}

\subsection{Column--Row Factorization and Rank of a Matrix}

We begin by defining two nonnegative integers associated with each matrix.

\begin{definition}[Column rank, row rank]\label{ax:3.52}
Suppose $A$ is an $m$-by-$n$ matrix with entries in $\F$.
\begin{itemize}
\item The \emph{column rank} of $A$ is the dimension of the span of the columns
  of $A$ in $\F^{m,1}$.
\item The \emph{row rank} of $A$ is the dimension of the span of the rows of
  $A$ in $\F^{1,n}$.
\end{itemize}
\end{definition}

If $A$ is an $m$-by-$n$ matrix, then the column rank of $A$ is at most $n$
(because $A$ has $n$ columns) and the column rank of $A$ is also at most $m$
(because $\dim\F^{m,1}=m$). Similarly, the row rank of $A$ is also at most
$\min\{m,n\}$.

\begin{example}[Column rank and row rank of a $2$-by-$4$ matrix]\label{ax:3.53}
Suppose
\[ A = \begin{pmatrix} 4 & 7 & 1 & 8\\ 3 & 5 & 2 & 9 \end{pmatrix}. \]

The column rank of $A$ is the dimension of
\[ \spn\biggl(\begin{pmatrix} 4\\ 3 \end{pmatrix},
   \begin{pmatrix} 7\\ 5 \end{pmatrix},
   \begin{pmatrix} 1\\ 2 \end{pmatrix},
   \begin{pmatrix} 8\\ 9 \end{pmatrix}\biggr) \]
in $\F^{2,1}$. Neither of the first two vectors listed above in $\F^{2,1}$ is a
scalar multiple of the other. Thus the span of this list of length four has
dimension at least two. The span of this list of vectors in $\F^{2,1}$ cannot
have dimension larger than two because $\dim\F^{2,1}=2$. Thus the span of this
list has dimension two, which means that the column rank of $A$ is two.

The row rank of $A$ is the dimension of
\[ \spn\Bigl(\begin{pmatrix} 4 & 7 & 1 & 8 \end{pmatrix},
   \begin{pmatrix} 3 & 5 & 2 & 9 \end{pmatrix}\Bigr) \]
in $\F^{1,4}$. Neither of the two vectors listed above in $\F^{1,4}$ is a
scalar multiple of the other. Thus the span of this list of length two has
dimension two, which means that the row rank of $A$ is two.
\end{example}

We now define the transpose of a matrix.

\begin{definition}[Transpose, $A^{\mathrm{t}}$]\label{ax:3.54}
The \emph{transpose} of a matrix $A$, denoted by $A^{\mathrm{t}}$, is the
matrix obtained from $A$ by interchanging rows and columns. Specifically, if
$A$ is an $m$-by-$n$ matrix, then $A^{\mathrm{t}}$ is the $n$-by-$m$ matrix
whose entries are given by the equation
\[ (A^{\mathrm{t}})_{k,j} = A_{j,k}. \]
\end{definition}

\begin{example}[Transpose of a matrix]\label{ax:3.55}
If $A=\begin{pmatrix} 5 & -7\\ 3 & 8\\ -4 & 2 \end{pmatrix}$, then
$A^{\mathrm{t}}=\begin{pmatrix} 5 & 3 & -4\\ -7 & 8 & 2 \end{pmatrix}$.

Note that here $A$ is a $3$-by-$2$ matrix and $A^{\mathrm{t}}$ is a
$2$-by-$3$ matrix.
\end{example}

The transpose has nice algebraic properties:
$(A+B)^{\mathrm{t}}=A^{\mathrm{t}}+B^{\mathrm{t}}$,
$(\lambda A)^{\mathrm{t}}=\lambda A^{\mathrm{t}}$, and
$(AC)^{\mathrm{t}}=C^{\mathrm{t}}A^{\mathrm{t}}$ for all $m$-by-$n$ matrices
$A,B$, all $\lambda\in\F$, and all $n$-by-$p$ matrices $C$ (see Exercises 14
and~15).

The next result will be the main tool used to prove that the column rank equals
the row rank (see \ref{ax:3.57}).

\begin{result}[Column--row factorization]\label{ax:3.56}
Suppose $A$ is an $m$-by-$n$ matrix with entries in $\F$ and column rank
$c\ge1$. Then there exist an $m$-by-$c$ matrix $C$ and a $c$-by-$n$ matrix $R$,
both with entries in $\F$, such that $A=CR$.
\end{result}

\begin{proof}
Each column of $A$ is an $m$-by-$1$ matrix. The list
$A_{\cdot,1},\dots,A_{\cdot,n}$ of columns of $A$ can be reduced to a basis of
the span of the columns of $A$ (by \ref{ax:2.30}). This basis has length $c$,
by the definition of the column rank. The $c$ columns in this basis can be put
together to form an $m$-by-$c$ matrix $C$.

If $k\in\{1,\dots,n\}$, then column $k$ of $A$ is a linear combination of the
columns of $C$. Make the coefficients of this linear combination into column
$k$ of a $c$-by-$n$ matrix that we call $R$. Then $A=CR$, as follows from
\ref{ax:3.51}(a).
\end{proof}

In Example~\ref{ax:3.53}, the column rank and row rank turned out to equal each
other. The next result states that this happens for all matrices.

\begin{result}[Column rank equals row rank]\label{ax:3.57}
Suppose $A\in\F^{m,n}$. Then the column rank of $A$ equals the row rank of $A$.
\end{result}

\begin{proof}
Let $c$ denote the column rank of $A$. If $c=0$, then $A=0$ and hence the row
rank of $A$ also equals $0$. Thus we can assume that $c\ge1$.

Let $A=CR$ be the column--row factorization of $A$ given by \ref{ax:3.56},
where $C$ is an $m$-by-$c$ matrix and $R$ is a $c$-by-$n$ matrix. Then
\ref{ax:3.51}(b) tells us that every row of $A$ is a linear combination of the
rows of $R$. Because $R$ has $c$ rows, this implies that the row rank of $A$ is
less than or equal to the column rank $c$ of $A$.

To prove the inequality in the other direction, apply the previous paragraph
result to $A^{\mathrm{t}}$, getting
\begin{align*}
\text{column rank of } A &= \text{row rank of } A^{\mathrm{t}}\\
  &\le \text{column rank of } A^{\mathrm{t}}\\
  &= \text{row rank of } A.
\end{align*}
Thus the column rank of $A$ equals the row rank of $A$.
\end{proof}

Because the column rank equals the row rank, the last result allows us to
dispense with the terms ``column rank'' and ``row rank'' and just use the
simpler term ``rank''.

\begin{definition}[Rank]\label{ax:3.58}
The \emph{rank} of a matrix $A\in\F^{m,n}$ is the column rank of $A$.
\end{definition}

See \ref{ax:3.133} and Exercise~8 in Section~7A for alternative proofs that the
column rank equals the row rank.

% ------------------------------------------------------------------ exercises
\begin{exc}

\begin{exercise}
Suppose $T\in\Lin(V,W)$. Show that with respect to each choice of bases of $V$
and $W$, the matrix of $T$ has at least $\dim\range T$ nonzero entries.
\end{exercise}
\begin{solution}
Let $v_1,\dots,v_n$ be a basis of $V$ and $w_1,\dots,w_m$ be a basis of $W$.
Let $A=\Mat(T)$ be the matrix with respect to these bases. Then the $j$-th
column of $A$ is the coordinate vector of $Tv_j$ with respect to the basis for
$W$.

Suppose $A$ had fewer than $\dim\range T$ nonzero columns. Since zero columns
do not contribute to the span, the columns of $A$ would then span a space of
dimension less than $\dim\range T$. This is a contradiction, since fewer than
$\dim\range T$ vectors cannot span $\range T$. Therefore, $A$ must have at
least $\dim\range T$ nonzero columns, and hence at least $\dim\range T$ nonzero
entries.
\end{solution}

\begin{exercise}
Suppose $T\in\Lin(V,W)$, where $V$ and $W$ are finite-dimensional and nonzero.
Prove that $\dim\range T=1$ if and only if there exist a basis of $V$ and a
basis of $W$ such that with respect to these bases, all entries of $\Mat(T)$
equal $1$.
\end{exercise}
\begin{solution}
Suppose $\dim\range T=1$. Then there exists a nonzero vector $y_1\in W$ such
that $\range T=\spn y_1$. Let $\dim V=n$ and $\dim W=m$. Extend $y_1$ to a
basis $y_1,\dots,y_m$ of $W$, and define the following vectors in $W$:
\[ w_1 = y_1-(y_2+\cdots+y_m), \quad w_k = y_k \text{ for } k=2,\dots,m. \]
It is clear that $w_1+w_2+\cdots+w_m=y_1$, and hence $w_1,\dots,w_m$ spans $W$.
Now suppose for scalars $\alpha_1,\dots,\alpha_m\in\F$, we have
\[ \alpha_1w_1+\alpha_2w_2+\cdots+\alpha_mw_m = 0. \]
Then we can express this as
\[ \alpha_1\bigl(y_1-(y_2+\cdots+y_m)\bigr)+\alpha_2y_2+\cdots+\alpha_my_m = 0. \]
Simplifying, we get
\[ (\alpha_1)y_1+(-\alpha_1+\alpha_2)y_2+\cdots+(-\alpha_1+\alpha_m)y_m = 0. \]
Since $y_1,\dots,y_m$ are linearly independent, we must have
$\alpha_1=\alpha_2=\cdots=\alpha_m=0$. Therefore, $w_1,\dots,w_m$ are linearly
independent and hence form a basis of $W$.

We now construct a basis of $V$. By the fundamental theorem of linear maps, we
have
\[ \dim V = \dim T+\dim\range T = \dim T+1. \]
Let $u_2,\dots,u_n$ be a basis of $T$. Then we can extend this to a basis
$u_1,u_2,\dots,u_n$ of $V$, where $u_1\notin T$. Since $\range T=\spn y_1$, we
have
\[ Tu_1 = \beta y_1 \]
for some nonzero scalar $\beta\in\F$. If $\beta\ne1$, we may replace $u_1$ with
$\beta^{-1}u_1$ to ensure that $Tu_1=y_1$. We construct the following vectors
in $V$:
\[ v_1 = u_1, \quad v_k = u_1+u_k \text{ for } k=2,\dots,n. \]
Note that for every $j=1,\dots,n$, we have $Tv_j=Tu_1+Tu_j=y_1+0=y_1$. We
claim that $v_1,\dots,v_n$ form a basis of $V$. Suppose for scalars
$\alpha_1,\dots,\alpha_n\in\F$, we have
\[ \alpha_1v_1+\alpha_2v_2+\cdots+\alpha_nv_n = 0. \]
Then we can express this as
\[ (\alpha_1+\alpha_2+\cdots+\alpha_n)u_1+\alpha_2u_2+\cdots+\alpha_nu_n = 0. \]
Since $u_1,\dots,u_n$ are linearly independent, we must have
$\alpha_1+\alpha_2+\cdots+\alpha_n=0$ and $\alpha_2=\cdots=\alpha_n=0$. This
implies that $\alpha_1=0$, and hence $v_1,\dots,v_n$ are linearly independent
and form a basis of $V$. Then with respect to the bases $v_1,\dots,v_n$ of $V$
and $w_1,\dots,w_m$ of $W$, we have that $Tv_j=y_1=w_1+w_2+\cdots+w_m$ for each
$j=1,\dots,n$. Therefore, the matrix of $T$ with respect to these bases has all
entries equal to $1$.

Conversely, suppose there exist a basis of $V$ and a basis of $W$ such that
with respect to these bases, all entries of $\Mat(T)$ equal $1$. Then for each
basis vector $v_j$ of $V$, we have
\[ Tv_j = w \]
for some fixed nonzero vector $w\in W$, where $w=w_1+\cdots+w_m$ is the sum of
the basis vectors of $W$. Since $\range T$ is spanned by the images of the basis
vectors of $V$, then any sum of $Tv_j$ is a scalar multiple of $w$. Therefore,
$\range T=\spn w$, and hence $\dim\range T=1$.
\end{solution}

\begin{exercise}
Suppose $v_1,\dots,v_n$ is a basis of $V$ and $w_1,\dots,w_m$ is a basis of
$W$.
\begin{parts}
\item Show that if $S,T\in\Lin(V,W)$, then $\Mat(S+T)=\Mat(S)+\Mat(T)$.
\item Show that if $\lambda\in\F$ and $T\in\Lin(V,W)$, then
  $\Mat(\lambda T)=\lambda\Mat(T)$.
\end{parts}
\exnote{This exercise asks you to verify \ref{ax:3.35} and \ref{ax:3.38}.}
\end{exercise}
\begin{solution}\leavevmode
\begin{parts}
\item Let $v_1,\dots,v_n$ be a basis of $V$ and $w_1,\dots,w_m$ be a basis of
  $W$. Then the matrix of $S+T$ with respect to these bases is given by
  \begin{align*}
  \Mat(S+T) &= \begin{pmatrix} (S+T)v_1 & (S+T)v_2 & \cdots & (S+T)v_n
               \end{pmatrix}\\
   &= \begin{pmatrix} Sv_1+Tv_1 & Sv_2+Tv_2 & \cdots & Sv_n+Tv_n
      \end{pmatrix}\\
   &= \begin{pmatrix} Sv_1 & Sv_2 & \cdots & Sv_n \end{pmatrix}
      + \begin{pmatrix} Tv_1 & Tv_2 & \cdots & Tv_n \end{pmatrix}\\
   &= \Mat(S)+\Mat(T).
  \end{align*}
\item Let $v_1,\dots,v_n$ be a basis of $V$ and $w_1,\dots,w_m$ be a basis of
  $W$. Then the matrix of $\lambda T$ with respect to these bases is given by
  \begin{align*}
  \Mat(\lambda T) &= \begin{pmatrix} (\lambda T)v_1 & (\lambda T)v_2 & \cdots
                     & (\lambda T)v_n \end{pmatrix}\\
   &= \begin{pmatrix} \lambda Tv_1 & \lambda Tv_2 & \cdots & \lambda Tv_n
      \end{pmatrix}\\
   &= \lambda\begin{pmatrix} Tv_1 & Tv_2 & \cdots & Tv_n \end{pmatrix}\\
   &= \lambda\Mat(T). \qedhere
  \end{align*}
\end{parts}
\end{solution}

\begin{exercise}
Suppose that $D\in\Lin\bigl(\Poly_3(\R),\Poly_2(\R)\bigr)$ is the
differentiation map defined by $Dp=p'$. Find a basis of $\Poly_3(\R)$ and a
basis of $\Poly_2(\R)$ such that the matrix of $D$ with respect to these bases
is
\[ \begin{pmatrix} 1 & 0 & 0 & 0\\ 0 & 1 & 0 & 0\\ 0 & 0 & 1 & 0
   \end{pmatrix}. \]
\exnote{Compare with Example~\ref{ax:3.33}. The next exercise generalizes this
exercise.}
\end{exercise}
\begin{solution}
Consider the standard basis of $\Poly_3(\R)$. Reorder it to obtain the basis
$x^3,x^2,x,1$. Consider the standard basis of $\Poly_2(\R)$. Reorder it to
obtain the basis $x^2,x,1$. Observe that $D(x^3)=3x^2$, $D(x^2)=2x$, $D(x)=1$,
and $D(1)=0$. However, we need the matrix of $D$ to have the desired form. To
achieve this, we can scale the basis vectors of $\Poly_3(\R)$ and $\Poly_2(\R)$
appropriately. Then we may define the basis for $\Poly_3(\R)$ as
\[ 1/3x^3,\,1/2x^2,\,x,\,1 \]
while maintaining the standard basis for $\Poly_2(\R)$ as $x^2,x,1$. Then we
have
\[ D\Bigl(1/3x^3\Bigr)=x^2, \quad D\Bigl(1/2x^2\Bigr)=x, \quad
   D(x)=1, \quad D(1)=0. \]
Therefore, with respect to the bases, we get the desired matrix representation
of $D$.
\end{solution}

\begin{exercise}
Suppose $V$ and $W$ are finite-dimensional and $T\in\Lin(V,W)$. Prove that
there exist a basis of $V$ and a basis of $W$ such that with respect to these
bases, all entries of $\Mat(T)$ are $0$ except that the entries in row $k$,
column $k$, equal $1$ if $1\le k\le\dim\range T$.
\end{exercise}
\begin{solution}
Let $n=\dim V$ and $m=\dim W$. By the fundamental theorem of linear maps, we
have
\[ \dim V = \dim\nul T+\dim\range T. \]
Let $r=\dim\range T$. Let $u_1,\dots,u_{n-r}$ be a basis of $\nul T$. Extend
this to a basis $v_1,\dots,v_r,u_1,\dots,u_{n-r}$ of $V$. We claim that
$Tv_1,\dots,Tv_r$ are linearly independent. Suppose for scalars
$\alpha_1,\dots,\alpha_r\in\F$, we have
\[ \alpha_1Tv_1+\cdots+\alpha_rTv_r = 0. \]
Then
\[ T(\alpha_1v_1+\cdots+\alpha_rv_r) = 0, \]
which implies that $\alpha_1v_1+\cdots+\alpha_rv_r\in\nul T$. Since
$v_1,\dots,v_r,u_1,\dots,u_{n-r}$ is a basis of $V$, it follows that
$\alpha_1=\cdots=\alpha_r=0$. Therefore, $Tv_1,\dots,Tv_r$ are linearly
independent and form a basis of $\range T$. Extend this to a basis
$Tv_1,\dots,Tv_r,w_1,\dots,w_{m-r}$ of $W$. Then with respect to these bases,
$Tv_k$ is the $k$-th basis vector of $W$ for each $k=1,\dots,r$, and
\[ Tu_j = 0, \quad\text{for } j=1,\dots,n-r. \]
Therefore, the matrix of $T$ with respect to these bases has $1$'s in the
diagonal entries corresponding to $Tv_k$ and $0$'s elsewhere, as desired.
\end{solution}

\begin{exercise}
Suppose $v_1,\dots,v_m$ is a basis of $V$ and $W$ is finite-dimensional.
Suppose $T\in\Lin(V,W)$. Prove that there exists a basis $w_1,\dots,w_n$ of $W$
such that all entries in the first column of $\Mat(T)$ [with respect to the
bases $v_1,\dots,v_m$ and $w_1,\dots,w_n$] are $0$ except for possibly a $1$ in
the first row, first column.
\exnote{In this exercise, unlike Exercise~5, you are given the basis of $V$
instead of being able to choose a basis of $V$.}
\end{exercise}
\begin{solution}
Let $n=\dim W$. If $Tv_1=0$, then we can choose any basis of $W$ and the first
column of $\Mat(T)$ will be all zeros. If $Tv_1\ne0$, then we can extend $Tv_1$
to a basis of $W$. Let $w_1=Tv_1$ and extend it to a basis
$w_1,w_2,\dots,w_n$ of $W$. Then with respect to the bases $v_1,\dots,v_m$ of
$V$ and $w_1,\dots,w_n$ of $W$, we have
\[ Tv_1 = w_1 \]
and for each $j=2,\dots,m$, we can express
\[ Tv_j = \alpha_{j1}w_1+\alpha_{j2}w_2+\cdots+\alpha_{jn}w_n \]
for some scalars $\alpha_{j1},\alpha_{j2},\dots,\alpha_{jn}\in\F$. Therefore,
the first column of $\Mat(T)$ has a $1$ in the first row and zeros in all other
rows, as desired.
\end{solution}

\begin{exercise}
Suppose $w_1,\dots,w_n$ is a basis of $W$ and $V$ is finite-dimensional.
Suppose $T\in\Lin(V,W)$. Prove that there exists a basis $v_1,\dots,v_m$ of $V$
such that all entries in the first row of $\Mat(T)$ [with respect to the bases
$v_1,\dots,v_m$ and $w_1,\dots,w_n$] are $0$ except for possibly a $1$ in the
first row, first column.
\exnote{In this exercise, unlike Exercise~5, you are given the basis of $W$
instead of being able to choose a basis of $W$.}
\end{exercise}
\begin{solution}
Let $m=\dim V$ and $u_1,\dots,u_m$ be a basis for $V$. If $T$ has trivial $w_1$
component, meaning that the coefficient of $w_1$ in $Tv$ is zero for all
$v\in V$, then choose any basis of $V$ and the first row of $\Mat(T)$ will be
all zeros. Otherwise, reorder the basis so that $Tu_1$ has non-trivial $w_1$
component. Then we can express
\[ Tu_1 = \alpha w_1+\beta_2w_2+\cdots+\beta_nw_n \]
for some scalars $\alpha,\beta_2,\dots,\beta_n\in\F$ such that $\alpha\ne0$. We
can scale $u_1$ by $\alpha^{-1}$ to obtain a new vector $v_1=\alpha^{-1}u_1$
such that
\[ Tv_1 = w_1+\frac{\beta_2}{\alpha}w_2+\cdots+\frac{\beta_n}{\alpha}w_n. \]
With the remaining $u_2,\dots,u_m$ basis vectors of $V$, note that they can be
expressed as
\[ Tu_j = \alpha_{j1}w_1+\alpha_{j2}w_2+\cdots+\alpha_{jn}w_n \]
for some scalars $\alpha_{j1},\alpha_{j2},\dots,\alpha_{jn}\in\F$. Consider the
set of vectors of $V$ given by
\[ v_1 = \alpha^{-1}u_1, \quad v_j = u_j-\alpha_{j1}v_1 \text{ for }
   j=2,\dots,m. \]
Note that $Tv_j=Tu_j-\alpha_{j1}Tv_1$ is just a combination of
$w_2,\dots,w_n$. We now claim that $v_1,v_2,\dots,v_m$ are linearly
independent. Suppose for scalars $\gamma_1,\dots,\gamma_m\in\F$, we have
\[ \gamma_1v_1+\gamma_2v_2+\cdots+\gamma_mv_m = 0. \]
Then we can express this as
\[ \gamma_1v_1+\gamma_2(u_2-\alpha_{21}v_1)+\cdots+\gamma_m(u_m-\alpha_{m1}v_1)
   = 0. \]
Simplifying, we get
\[ (\gamma_1-\gamma_2\alpha_{21}-\cdots-\gamma_m\alpha_{m1})v_1
   +\gamma_2u_2+\cdots+\gamma_mu_m = 0. \]
Since $v_1,u_2,\dots,u_m$ are linearly independent, we must have
$\gamma_1-\gamma_2\alpha_{21}-\cdots-\gamma_m\alpha_{m1}=0$ and
$\gamma_2=\cdots=\gamma_m=0$. This implies that $\gamma_1=0$, and hence
$v_1,v_2,\dots,v_m$ are linearly independent. Then $v_1,v_2,\dots,v_m$ form a
basis of $V$. It follows that $Tv_1$ has a $1$ in the first row, while $Tv_j$
for $j=2,\dots,m$ have zeros in the first row.
\end{solution}

\begin{exercise}
Suppose $A$ is an $m$-by-$n$ matrix and $B$ is an $n$-by-$p$ matrix. Prove
that
\[ (AB)_{j,\cdot} = A_{j,\cdot}B \]
for each $1\le j\le m$. In other words, show that row $j$ of $AB$ equals (row
$j$ of $A$) times $B$.
\exnote{This exercise gives the row version of \ref{ax:3.48}.}
\end{exercise}
\begin{solution}
Let $A$ be an $m$-by-$n$ matrix and $B$ be an $n$-by-$p$ matrix. The entry in
row $j$, column $k$, of the product $AB$ is given by
\[ (AB)_{j,k} = \sum_{r=1}^n A_{j,r}B_{r,k}. \]
The row vector $A_{j,\cdot}$ is the $j$-th row of $A$, which can be expressed
as
\[ A_{j,\cdot} = (A_{j,1},A_{j,2},\dots,A_{j,n}). \]
When we multiply this row vector by the matrix $B$, we get a new row vector
whose entries are given by
\[ (A_{j,\cdot}B)_k = \sum_{r=1}^n A_{j,r}B_{r,k}. \]
Comparing this with the expression for $(AB)_{j,k}$, we see that
\[ (AB)_{j,k} = (A_{j,\cdot}B)_k \]
for each column index $k$. Therefore, the entire row vector $(AB)_{j,\cdot}$ is
equal to the row vector $A_{j,\cdot}B$. This proves that row $j$ of $AB$ equals
(row $j$ of $A$) times $B$.
\end{solution}

\begin{exercise}
Suppose $a=\begin{pmatrix} a_1 & \cdots & a_n \end{pmatrix}$ is a $1$-by-$n$
matrix and $B$ is an $n$-by-$p$ matrix. Prove that
\[ aB = a_1B_{1,\cdot}+\cdots+a_nB_{n,\cdot}. \]
In other words, show that $aB$ is a linear combination of the rows of $B$, with
the scalars that multiply the rows coming from $a$.
\exnote{This exercise gives the row version of \ref{ax:3.50}.}
\end{exercise}
\begin{solution}
Let $a=\begin{pmatrix} a_1 & \cdots & a_n \end{pmatrix}$ be a $1$-by-$n$
matrix and $B$ be an $n$-by-$p$ matrix. The product $aB$ is a $1$-by-$p$
matrix, and its entries are given by
\[ (aB)_{1,k} = \sum_{j=1}^n a_jB_{j,k} \]
for each column index $k$.

Now, consider the linear combination of the rows of $B$ given by
\[ a_1B_{1,\cdot}+a_2B_{2,\cdot}+\cdots+a_nB_{n,\cdot}. \]
The entry in row $1$, column $k$, of this linear combination is
\begin{multline*}
(a_1B_{1,\cdot}+a_2B_{2,\cdot}+\cdots+a_nB_{n,\cdot})_k\\
  = a_1B_{1,k}+a_2B_{2,k}+\cdots+a_nB_{n,k} = \sum_{j=1}^n a_jB_{j,k}.
\end{multline*}

Comparing this with the expression for $(aB)_{1,k}$, we see that
\[ (aB)_{1,k} = (a_1B_{1,\cdot}+a_2B_{2,\cdot}+\cdots+a_nB_{n,\cdot})_k \]
for each column index $k$. Therefore, we conclude that
\[ aB = a_1B_{1,\cdot}+a_2B_{2,\cdot}+\cdots+a_nB_{n,\cdot}, \]
which shows that $aB$ is indeed a linear combination of the rows of $B$, with
the scalars coming from the entries of $a$.
\end{solution}

\begin{exercise}
Give an example of $2$-by-$2$ matrices $A$ and $B$ such that $AB\ne BA$.
\end{exercise}
\begin{solution}
Let
\[ A = \begin{pmatrix} 1 & 2\\ 3 & 4 \end{pmatrix} \quad\text{and}\quad
   B = \begin{pmatrix} 0 & 1\\ 1 & 0 \end{pmatrix}. \]
Then we compute the products $AB$ and $BA$:
\[ AB = \begin{pmatrix} 1 & 2\\ 3 & 4 \end{pmatrix}
        \begin{pmatrix} 0 & 1\\ 1 & 0 \end{pmatrix}
      = \begin{pmatrix} 1\cdot0+2\cdot1 & 1\cdot1+2\cdot0\\
                        3\cdot0+4\cdot1 & 3\cdot1+4\cdot0 \end{pmatrix}
      = \begin{pmatrix} 2 & 1\\ 4 & 3 \end{pmatrix}, \]
and
\[ BA = \begin{pmatrix} 0 & 1\\ 1 & 0 \end{pmatrix}
        \begin{pmatrix} 1 & 2\\ 3 & 4 \end{pmatrix}
      = \begin{pmatrix} 0\cdot1+1\cdot3 & 0\cdot2+1\cdot4\\
                        1\cdot1+0\cdot3 & 1\cdot2+0\cdot4 \end{pmatrix}
      = \begin{pmatrix} 3 & 4\\ 1 & 2 \end{pmatrix}. \]
By comparing the two products, we see that $AB\ne BA$.
\end{solution}

\begin{exercise}
Prove that the distributive property holds for matrix addition and matrix
multiplication. In other words, suppose $A$, $B$, $C$, $D$, $E$, and $F$ are
matrices whose sizes are such that $A(B+C)$ and $(D+E)F$ make sense. Explain
why $AB+AC$ and $DF+EF$ both make sense and prove that
\[ A(B+C) = AB+AC \quad\text{and}\quad (D+E)F = DF+EF. \]
\end{exercise}
\begin{solution}
Let $A$ be an $m$-by-$n$ matrix, and let $B$ and $C$ be $n$-by-$p$ matrices.
Then $B+C$ is an $n$-by-$p$ matrix, and hence $A(B+C)$ is an $m$-by-$p$
matrix. The products $AB$ and $AC$ are also defined, and both are $m$-by-$p$
matrices. Therefore, the sum $AB+AC$ is defined and is also an $m$-by-$p$
matrix. By Exercise 3C8 and Exercise 3C9, we have
\begin{multline*}
A(B+C)_{j,\cdot} = A_{j,\cdot}(B+C) = A_{j,\cdot}B+A_{j,\cdot}C\\
  = (AB)_{j,\cdot}+(AC)_{j,\cdot} = (AB+AC)_{j,\cdot}
\end{multline*}
for all $j=1,\dots,m$. Therefore, $A(B+C)=AB+AC$.

Similarly, let $D$ and $E$ be $m$-by-$n$ matrices, and let $F$ be an $n$-by-$p$
matrix. Then $D+E$ is an $m$-by-$n$ matrix, and hence $(D+E)F$ is an
$m$-by-$p$ matrix. The products $DF$ and $EF$ are also defined, and both are
$m$-by-$p$ matrices. Therefore, the sum $DF+EF$ is defined and is also an
$m$-by-$p$ matrix. Then the product $(D+E)F$ is defined. By Exercise 3C8 and
Exercise 3C9, we have
\begin{multline*}
(D+E)F_{j,\cdot} = (D+E)_{j,\cdot}F = D_{j,\cdot}F+E_{j,\cdot}F\\
  = (DF)_{j,\cdot}+(EF)_{j,\cdot} = (DF+EF)_{j,\cdot}
\end{multline*}
for all $j=1,\dots,m$. Therefore, $(D+E)F=DF+EF$.
\end{solution}

\begin{exercise}
Prove that matrix multiplication is associative. In other words, suppose $A$,
$B$, and $C$ are matrices whose sizes are such that $(AB)C$ makes sense.
Explain why $A(BC)$ makes sense and prove that
\[ (AB)C = A(BC). \]
\exnote{Try to find a clean proof that illustrates the following quote from Emil
Artin: ``It is my experience that proofs involving matrices can be shortened
by 50\% if one throws the matrices out.''}
\end{exercise}
\begin{solution}
\solnote{Manual Remark}{Unfortunately, we cannot use the same trick as in
Exercise 3C11 to prove this result.}

Let $A$ be an $m$-by-$n$ matrix, $B$ be an $n$-by-$p$ matrix, and $C$ be a
$p$-by-$q$ matrix. Then the product $AB$ is defined and is an $m$-by-$p$
matrix. Therefore, the product $(AB)C$ is also defined and is an $m$-by-$q$
matrix. Going by the definition of matrix multiplication, we have
\[ \bigl((AB)C\bigr)_{j,k} = \sum_{r=1}^p (AB)_{j,r}C_{r,k}
   = \sum_{r=1}^p\Bigl(\sum_{s=1}^n A_{j,s}B_{s,r}\Bigr)C_{r,k}
   = \sum_{r=1}^p\sum_{s=1}^n A_{j,s}B_{s,r}C_{r,k}. \]
On the other hand, the product $BC$ is defined and is an $n$-by-$q$ matrix.
Therefore, the product $A(BC)$ is also defined and is an $m$-by-$q$ matrix.
Again, by the definition of matrix multiplication, we have
\[ \bigl(A(BC)\bigr)_{j,k} = \sum_{s=1}^n A_{j,s}(BC)_{s,k}
   = \sum_{s=1}^n A_{j,s}\Bigl(\sum_{r=1}^p B_{s,r}C_{r,k}\Bigr)
   = \sum_{s=1}^n\sum_{r=1}^p A_{j,s}B_{s,r}C_{r,k}. \]
Since the order of summation does not matter, we can interchange the order of
the sums in the last expression to get
\[ \bigl(A(BC)\bigr)_{j,k} = \sum_{r=1}^p\sum_{s=1}^n A_{j,s}B_{s,r}C_{r,k}. \]
Comparing the expressions for $\bigl((AB)C\bigr)_{j,k}$ and
$\bigl(A(BC)\bigr)_{j,k}$, we see that they are equal for all $j$ and $k$.
Therefore, $(AB)C=A(BC)$.
\end{solution}

\begin{exercise}
Suppose $A$ is an $n$-by-$n$ matrix and $1\le j,k\le n$. Show that the entry
in row $j$, column $k$, of $A^3$ (which is defined to mean $AAA$) is
\[ \sum_{p=1}^n\sum_{r=1}^n A_{j,p}A_{p,r}A_{r,k}. \]
\end{exercise}
\begin{solution}
Note that $A^3=(AA)A=A^2A$. By the definition of matrix multiplication, we have
\[ (A^3)_{j,k} = (A^2A)_{j,k} = \sum_{p=1}^n (A^2)_{j,p}A_{p,k}. \]
Now, we can express $(A^2)_{j,p}$ as
\[ (A^2)_{j,p} = (AA)_{j,p} = \sum_{r=1}^n A_{j,r}A_{r,p}. \]
Substituting this back into the expression for $(A^3)_{j,k}$, we get
\[ (A^3)_{j,k} = \sum_{p=1}^n\Bigl(\sum_{r=1}^n A_{j,r}A_{r,p}\Bigr)A_{p,k}
   = \sum_{p=1}^n\sum_{r=1}^n A_{j,r}A_{r,p}A_{p,k}. \]
By renaming the indices $r$ and $p$, we can rewrite this as
\[ (A^3)_{j,k} = \sum_{p=1}^n\sum_{r=1}^n A_{j,p}A_{p,r}A_{r,k}, \]
which is the desired expression.
\end{solution}

\begin{exercise}
Suppose $m$ and $n$ are positive integers. Prove that the function
$A\mapsto A^{\mathrm{t}}$ is a linear map from $\F^{m,n}$ to $\F^{n,m}$.
\end{exercise}
\begin{solution}
Let $A,B\in\F^{m,n}$ and let $\lambda\in\F$. We need to show that
$(A+B)^t=A^t+B^t$ and $(\lambda A)^t=\lambda A^t$.

First, consider the sum $A+B$. The entry in row $j$, column $k$, of $A+B$ is
given by
\[ (A+B)_{j,k} = A_{j,k}+B_{j,k}. \]
Taking the transpose, we have
\[ \bigl((A+B)^t\bigr)_{k,j} = (A+B)_{j,k} = A_{j,k}+B_{j,k}
   = (A^t)_{k,j}+(B^t)_{k,j} = (A^t+B^t)_{k,j}. \]
Therefore, $(A+B)^t=A^t+B^t$.

Next, consider the scalar multiplication $\lambda A$. The entry in row $j$,
column $k$, of $\lambda A$ is given by
\[ (\lambda A)_{j,k} = \lambda A_{j,k}. \]
Taking the transpose, we have
\[ \bigl((\lambda A)^t\bigr)_{k,j} = (\lambda A)_{j,k} = \lambda A_{j,k}
   = \lambda(A^t)_{k,j} = (\lambda A^t)_{k,j}. \]
Therefore, $(\lambda A)^t=\lambda A^t$.

Since both properties hold, we conclude that the function $A\mapsto A^t$ is a
linear map from $\F^{m,n}$ to $\F^{n,m}$.
\end{solution}

\begin{exercise}
Prove that if $A$ is an $m$-by-$n$ matrix and $C$ is an $n$-by-$p$ matrix,
then
\[ (AC)^{\mathrm{t}} = C^{\mathrm{t}}A^{\mathrm{t}}. \]
\exnote{This exercise shows that the transpose of the product of two matrices is
the product of the transposes in the opposite order.}
\end{exercise}
\begin{solution}
Let $A$ be an $m$-by-$n$ matrix and $C$ be an $n$-by-$p$ matrix. The entry in
row $j$, column $k$, of the product $AC$ is given by
\[ (AC)_{j,k} = \sum_{r=1}^n A_{j,r}C_{r,k}. \]
Taking the transpose, we have
\[ \bigl((AC)^t\bigr)_{k,j} = (AC)_{j,k} = \sum_{r=1}^n A_{j,r}C_{r,k}. \]

Now, consider the product $C^tA^t$. The entry in row $k$, column $j$, of this
product is given by
\[ (C^tA^t)_{k,j} = \sum_{r=1}^n (C^t)_{k,r}(A^t)_{r,j}
   = \sum_{r=1}^n C_{r,k}A_{j,r}. \]

Since multiplication in the field $\F$ is commutative, we can rewrite this as
\[ (C^tA^t)_{k,j} = \sum_{r=1}^n A_{j,r}C_{r,k}, \]
which is exactly the same as the expression for $\bigl((AC)^t\bigr)_{k,j}$.

Therefore, we conclude that
\[ (AC)^t = C^tA^t. \qedhere \]
\end{solution}

\begin{exercise}
Suppose $A$ is an $m$-by-$n$ matrix with $A\ne0$. Prove that the rank of $A$
is $1$ if and only if there exist $(c_1,\dots,c_m)\in\F^m$ and
$(d_1,\dots,d_n)\in\F^n$ such that $A_{j,k}=c_jd_k$ for every $j=1,\dots,m$
and every $k=1,\dots,n$.
\end{exercise}
\begin{solution}
Suppose the rank of $A$ is $1$. Then the column rank of $A$ is $1$, which means
that all columns of $A$ are scalar multiples of a single non-zero column
vector. Let $c=(c_1,\dots,c_m)^t$ be this non-zero column vector. Then for each
column $k$ of $A$, there exists a scalar $d_k\in\F$ such that the $k$-th column
of $A$ is given by $d_kc$. Therefore, we can express the entries of $A$ as
\[ A_{j,k} = c_jd_k \]
for all $j=1,\dots,m$ and $k=1,\dots,n$.

Conversely, suppose there exist $(c_1,\dots,c_m)\in\F^m$ and
$(d_1,\dots,d_n)\in\F^n$ such that $A_{j,k}=c_jd_k$ for all $j$ and $k$. Since
$A\ne0$, then $(c_1,\dots,c_m)\ne0$ and $(d_1,\dots,d_n)\ne0$. This means that
all columns of $A$ are scalar multiples of the column vector $c$, and hence the
column rank of $A$ is $1$. Therefore, the rank of $A$ is $1$.
\end{solution}

\begin{exercise}
Suppose $T\in\Lin(V)$, and $u_1,\dots,u_n$ and $v_1,\dots,v_n$ are bases of
$V$. Prove that the following are equivalent.
\begin{parts}
\item $T$ is injective.
\item The columns of $\Mat(T)$ are linearly independent in $\F^{n,1}$.
\item The columns of $\Mat(T)$ span $\F^{n,1}$.
\item The rows of $\Mat(T)$ span $\F^{1,n}$.
\item The rows of $\Mat(T)$ are linearly independent in $\F^{1,n}$.
\end{parts}
Here $\Mat(T)$ means $\Mat\bigl(T,(u_1,\dots,u_n),(v_1,\dots,v_n)\bigr)$.
\end{exercise}
\begin{solution}\leavevmode
\begin{itemize}
\item (a) $\iff$ (b): (a) $\iff$ $\dim\nul T=0$ $\iff$ $\dim\range T=n$
  $\iff$ $\operatorname{column\,rank}\Mat(T)=n$. For $n$ column vectors,
  rank $n$ implies linear independence $\iff$ (b).
\item (b) $\iff$ (c): $n$ columns of $\Mat(T)$ are linearly independent in
  $\F^{n,1}$ $\iff$ $n$ columns of $\Mat(T)$ span $\F^{n,1}$.
\item (c) $\iff$ (d): columns span $\F^{n,1}$ $\iff$ $\rank\Mat(T)=n$ by
  (b) $\iff$ (c). By Theorem~\ref{ax:3.57},
  $\rank\Mat(T)=\operatorname{row\,rank}\Mat(T)$ $\iff$ rows span $\F^{1,n}$.
\item (d) $\iff$ (e): $n$ rows of $\Mat(T)$ span $\F^{1,n}$ $\iff$ $n$ rows of
  $\Mat(T)$ are linearly independent in $\F^{1,n}$.
\end{itemize}
\end{solution}
\end{exc}
